% section: 等比数列の総和とその導出

\section{等比数列の和}\label{節：等比数列の和}
\sectionref*{節：階差型}と\sectionref*{節：特性方程式型}の別解で使う.
初項$a$,公比$r$の\bookterm{geometric}{等比数列}の和を$S_n=\sum_{k=1}^n ar^{k-1}$とする.

\subsection{和の公式}\label{節：等比数列の和：和の公式}
\[
  S_n=\frac{a(1-r^n)}{1-r}=\frac{a(r^n-1)}{r-1}
  \quad(r\ne1).
\]
$r=1$では全項が$a$なので$S_n=na$である.

\subsection{公式の導出}\label{節：等比数列の和：公式の導出}
総和を$r$倍すると,元の総和の項と位置がそろう.引けば途中の項が消える.
\BookMathLayout{\[
    \begin{array}{rrrl}
      &S_n={}&a+&\hspace{-0.75em}ar+ar^2+\cdots+ar^{n-1}\\
      -\,) &rS_n={}&{}&\hspace{-0.75em}ar+ar^2+ar^3+\cdots+ar^n\\ \hline
      &(1-r)S_n={}&{}&\hspace{-0.75em}a-ar^n
    \end{array}
  \]}{
  \[
    \begin{array}{rl}
      S_n={}&a+ar+ar^2+\cdots\\
      &{}+ar^{n-1}\\
      -\,)\ rS_n={}&ar+ar^2+ar^3+\cdots\\
      &{}+ar^n\\ \hline
      (1-r)S_n={}&a-ar^n
    \end{array}
  \]
}
$r\ne1$なら$1-r$で割り,公式を得る.倍率をそろえて差を取る工夫である.
